\documentclass[12pt]{article}
\usepackage[a4paper, total={6.5in, 9in}]{geometry}

%%%%%%%%%%%%%%%%%%%%%%%%
% Import packages
%%%%%%%%%%%%%%%%%%%%%%%%

\usepackage[a4paper, total={6in, 8in}]{geometry}
\usepackage{graphicx} % Required for inserting images
\usepackage{amsmath} % Gives nice maths
\usepackage{amssymb} % extra maths symbols
\usepackage{gensymb} % extra maths symbols
\usepackage{enumerate} % Gives better listing options including roman numberal
\usepackage{listings} % Used for code blocks (listings)
\usepackage{xcolor}
\usepackage{braket}
\usepackage{float}

\usepackage{subcaption} % used for subfigures

\usepackage[symbol]{footmisc} % TODO figure out a better footnoting scheme - using numbers is confusing but I'm not sure I like this package's default either

\usepackage{parskip} % Parskip means paragraphs get a gap between them instead of an indented start

%%%%%%%%%%%%%%%%%%%%%%%%%
% Aliases for common patterns
%%%%%%%%%%%%%%%%%%%%%%%%%

\newcommand{\secs}{\text{ s}} % Shorthand for 'seconds' unit in maths mode

\newcommand{\rr}[1]{\frac{1}{\sqrt{#1}}} % Reciprocal root e.g. \rr{2}
\newcommand{\boxpar}[1]{\fbox{\parbox{\textwidth}{#1}}} % Boxed paragraph

\newcommand{\mb}[1]{\begin{matrix}#1\end{matrix}}
\newcommand{\m}[1]{\left[\mb{#1}\right]}

\newcommand{\br}[1]{\left(#1\right)}

\newcommand{\rad}{\text{ rad}}

\newcommand{\questeq}{\overset{?}{=}}

\newcommand{\comment}[1]{}

%%%%%%%%%%%%%%%%%%%%%%%%
% Configuration for code listings
%%%%%%%%%%%%%%%%%%%%%%%%

% Define colors for code listings
\definecolor{codegreen}{rgb}{0,0.6,0}
\definecolor{codegray}{rgb}{0.5,0.5,0.5}
\definecolor{codepurple}{rgb}{0.58,0,0.82}
\definecolor{backcolour}{rgb}{0.95,0.95,0.92}

% Define a style for displaying assembly code
\lstdefinestyle{asm}{
  language=[x86masm]Assembler,
  basicstyle=\ttfamily\small,
  numbers=left,
  numberstyle=\tiny\color{gray},
  stepnumber=1,
  numbersep=5pt,
  backgroundcolor=\color{white},
  showspaces=false,
  showstringspaces=false,
  showtabs=false,
  frame=single,
  rulecolor=\color{black},
  tabsize=4,
  captionpos=b,
  breaklines=true,
  breakatwhitespace=false,
  keywordstyle=\color{blue},
  commentstyle=\color{green!50!black},
  stringstyle=\color{red},
  morekeywords={add, addi, sub, and, or, xor, sll, srl, sra, lui, auipc, 
          lb, lh, lw, lbu, lhu, sb, sh, sw, beq, bne, blt, bge, 
          bltu, bgeu, jal, jalr, ecall, ebreak}, % Add RISC-V specific keywords
    basicstyle=\ttfamily\small,
    \singlespacing
}
\lstset{style=mystyle} % "mystyle" code listing set
 % Import configuration (packages, styles, defaults, etc.)
\usepackage[defernumbers=true]{biblatex}

% Create a 'cited' category that gets automatically applied to everything cited in the text
\DeclareBibliographyCategory{cited}
\AtEveryCitekey{\addtocategory{cited}{\thefield{entrykey}}}

\nocite{*}

% anything in the .bib file with the 'cited' category will appear in the references
% anything in the .bib not cited in the text will appear in the bibliography

% add the bib file:
\addbibresource{bibliography.bib} % Import bibliography configuration

\begin{document}

%%%%%%%%%%%%%%%%%%%%%%%%%%%%%%%%%%%%%
% Cover page and table of contents
%%%%%%%%%%%%%%%%%%%%%%%%%%%%%%%%%%%%%
\begin{figure}
    \centering
    
\includegraphics[width=0.25\linewidth]{ucdlogo.png}
\end{figure}
\begin{centering}
    Electronic \& Computer Engineering\\
    EEEN40280 Lab 1\\
    2025/26 \\[20pt]
    Accurate frqeuency synthesis on ADuC841\\
    Thomas O'Dea (21386441)\\
    Aidan Quigney (XXXXXXXX)\\
    Tinu's full name? (XXXXXXXXX)\\[20pt]
\end{centering}

\textbf{Declaration of Authorship}\\
We declare that all material in this assessment is our own work except where there is clear acknowledgement and appropriate reference to the work of others.

\today

\pagebreak

\tableofcontents
\listoffigures
\listoftables

\pagebreak
% TODO maybe remove list of figures/tables if we don't have many

%%%%%%%%%%%%%%%%%%%%%%%%%%%%%%%%%%%%%
% Report body
%%%%%%%%%%%%%%%%%%%%%%%%%%%%%%%%%%%%%

\section*{Introduction}

The aim of this laboratory exercise was to design and implement a frequency synthesiser using an ADuC841 microcontroller programmed in assembly language. The lab was intended to build understanding of timing and frequency generation in embedded systems, starting with software-based methods and progressing to hardware based solutions using timers and interrupts. 

As an initial step, we used the example program blink.asm as a reference, this program demonstrated the basic assembly language structure, input/output pin control and the use of delay loops.

In the first part of the assignment, a square-wave was generated using software delays only, without the use of hardware timers or interrupts. The output frequency was controlled by calculating the instruction execution times and then adjusting the delay loops accordingly. 

In t

- Lab was to create a frequency synthesiser
- blink.asm as a precursor
- first in software, then in hardware

\section*{Task 1: Software Timing}
- show blink.asm as a start point then our changes
- show calculations for runtime of DELAY function (I have a spreadsheet) and calculate desired number of clock cycles to synthesise our target frequency
- show screenshot of waveform
- highlight unreliable timing in a software implementation which could cause frequency jitters

\subsection*{Program Design}

A sample program \texttt{blink.asm} as provided for initial demonstration of using software delays to generate a frequency. The target frequency was much lower and used to drive an LED rather than a transducer, but the software is similar.

This program was modified to provide a frequency of 900 Hz using a nested for loop to delay the CPU core. The precise number of clock cycles to generate a 900 Hz square wave was calculated and the delay function tweaked accordingly. Only two levels of 'for loop' nesting was necessary for this faster period, rather than three levels in the original \texttt{blink.asm}.

The main program loop has only three operations:\begin{itemize}
    \item Call a 'delay' to stall the CPU for a time
    \item Complement (flip) the output bit labelled 'SOUND' which is tied to the transducer on the physical board
    \item Return to the start of the loop
\end{itemize}

The ADuC841 runs with a clock speed of 11.0592 MHz. To accomplish a 900 Hz wave, the precise duration of this loop had to be precisely 6144 Hz, as given by equation \ref{eq:sw-timing-1}. Note the output bit must flip twice per wave period.

\begin{equation}
    \label{eq:sw-timing-1}
    11.0592 \times 10^6 \frac{\text{clock cycles}}{\text{s}}=6144 \frac{\text{clock cycles}}{\text{bit complement}}\times (2\times900) \frac{\text{bit complements}}{\text{s}}
\end{equation}

A delay function was designed to generate this precise period, accounting for the overhead in each loop. A table calculating the precise duration of the loop is shown in table \ref{tab:sw-timing-loop-duration}. Two \texttt{NOP} operations were included to artificially increase the loop duration.


\begin{table}
    \centering
    \begin{tabular}{lllllll}
        & Instruction & Clock Cycles & Occurences per loop\\
        \hline
        MAIN: & & &\\
        LOOP: & CALL DELAY    & 3+delay & 1\\
        &CPL SOUND            & 2 & 1\\
        &JMP LOOP             & 3 & 1\\
        \hline
        DELAY: &MOV R5, \#16  & 2 & 1\\
        DLY1: & MOV R7, \#126 & 2 & 16\\
              &DJNZ R7, \$    & 3 & $2016 = 16\times 126$\\
              &DJNZ R5, DLY1  & 3 & 16\\
              &NOP            & 1 & 1\\
              &NOP            & 1 & 1\\
              &RET            & 4 & 1\\
        \hline
        \textbf{Total} && 6144
    \end{tabular}
    \caption{Calculation of main loop duration in clock cycles. See the actual Assembly in listing \ref{lst:software-timing}.}
    \label{tab:sw-timing-loop-duration}
\end{table}

\begin{table}
    \centering
    \begin{tabular}{lllllll}
        & Instruction & Clock Cycles & Total occurences per loop\\
        \hline
        LOOP: & CALL DELAY    & 3 & 1\\
        DELAY: &MOV R5, \#16   & 2 & 1\\
        DLY1: & MOV R7, \#126 & 2 & 16\\
        &DJNZ R7, \$          & 3 & $2016 = 16\times 126$\\
        &DJNZ R5, DLY1        & 3 & 16\\
        &NOP                  & 1 & 1\\
        &NOP                  & 1 & 1\\
        &RET                  & 4 & 1\\
        &CPL SOUND            & 2 & 1\\
        &JMP LOOP             & 3 & 1\\
        \hline
        \textbf{Total} && 6144
    \end{tabular}
    \caption{Calculation of main loop duration in clock cycles. Note some code has been inlined for readability - compare with the actual Assembly in listing \ref{lst:software-timing}.}
    \label{tab:sw-timing-loop-duration}
\end{table}

\subsection*{Results}
TODO:

- Show simulation waveform with precise period as expected

- Show oscilloscope measurements, both time domain and FFT

- Show photograph of board setup

\subsubsection*{Simulation}
The code was first simulated in Keil uVision to verify performance. The number of clock cycles was calculated as follows:\begin{itemize}
    \item Place a breakpoint on the command "CPL SOUND"
    \item Run the program to this point
    \item Note the "states" count for the simulation, indicating how many clock cycles have run so far
    \item Run the program until the breakpoint arrives again
    \item Note the "state" counts again and subtract the previous count, giving the total clock cycles from one bit flip to the next
\end{itemize}

The period was also measured using the logic analyser which gave a result in simulation time rather than clock cycles. This is shown in figure \ref{TODO}. All measurements were consistent with expectations.

TODO screenshot of uVision simulation(s)

\subsubsection*{Physical Measurement}

The appropriate port on the ADuC841 was disconnected from the transducer and instead connected to the oscilloscope for measuring the square wave output. Figure \ref{fig:sw-timing-oscilloscope-measurements} shows both a time-domain and frequency-domain measurements.

\begin{figure}
    \centering
    \includegraphics[width=0.49\linewidth]{photos/900Hz.PNG}
    \includegraphics[width=0.49\linewidth]{photos/software_fft.PNG}
    \caption{Oscilloscope measurements of the 900 Hz square wave synthesised with software in both time domain (left) and frequency domain (right).}
    \label{fig:sw-timing-oscilloscope-measurements}
\end{figure}

The oscilloscope measurement in time domain shows a frequency of 900.22 Hz, an error of only 244 parts-per-million. This error may arise from imperfections in the frequency of the 11.0592 MHz clock, or from the oscilloscope measurement itself.

The FFT view shows a fundamental frequency of 895.18 Hz. This is simply due to the granularity of the FFT jumping 20 Hz at a time. The expected series of harmonics for a square wave is clearly present at odd multiples of the fundamental frequencies - the first such harmonic is highlighted with a cursor at a frequency of 2.71 kHz and an amplitude one-third of the fundamental's. Small spurs are present around each of the harmonics due to imperfections in the square wave.

\subsection*{Code}

The full program is shown in listing \ref{lst:software-timing}.

\lstinputlisting[caption=Source code for generating 900 Hz wave purely in software,label=lst:software-timing]{software_frequency/Source/software_frequency.a51}
% TODO template says to include code but it's very long and bulky. We can slim down the comment and reduce font size but still. Could move to an appendix



\section*{Example Program}
- We predicted a frequency of 675 Hz (show spreadsheet with calculations)
- We measured XYZ? (show oscilloscope screenshot)

\section*{Task 2: Hardware timing}
- discuss three options suggested by the handout
- show calculations: using timer 2 with a reload mechanism provides frequencies from 169 Hz to 11 MHz (or half that? check), with the target range of 600 Hz - 9 kHz having fine enough granularity to ensure a frequency error under 0.07\%
- discuss decision to rely on polling vs. interrupts for:
    - frequency synthesis
    - heartbeat LED
    - user button input
    - user switches input + output
- note that frequency synthesis ISR has priority to minimise any frequency jitter that could otherwise be caused by other interrupts
- list general purpose registers and how they are used by this program
- discuss use of lookup tables instead of calculations. Show how DPTR and A are used to access LUTs
- show formula for setting values in LUTs for a given frequency
- discuss technique for debouncing (and alternatives that were considered e.g. using TIC instead. If TIC were enabled we could also disable timer 1 and rely on the main loop for the heartbeat)
- note steps taken to ensure the heartbeat LED updates immediately when the user presses the button
- discuss resource utilisation (e.g. two timers, an external interrupt, etc.) and how this could be reduced if necessary

- attach code
- attach screenshot of waveforms
- attach working video?

\subsection*{Program Design}
\subsection*{Improvement}
\subsection*{Results}


\section*{Conclusion}


%%%%%%%%%%%%%%%%%%%%%%%%%%%%%%%%%%%%%
% References and bibliography and appendix
%%%%%%%%%%%%%%%%%%%%%%%%%%%%%%%%%%%%%
\pagebreak

% TODO we could remove these if we don't have any references to put
\section{References}
% A list of references formatted in a consistent manner.
\printbibliography[category=cited,heading=none]

\section{Bibliography}
% A bibliography (distinct from the references section) should be inlcuded
\printbibliography[notcategory=cited,heading=none]

\end{document}